\section{Introduction}
Next-scale representations describe an image as complete coarse maps followed by residual refinements at finer scales~\cite{var}. Each scale spans the full latent field and conditions the next finer scale, yielding a progressively refinable description and a generative prior for untransmitted content.

Under a fixed per-image symbol budget, transmitting only complete scale maps leaves capacity unused whenever the next residual map exceeds the remaining capacity, even if a prefix of its tokens would fit.

ARPC selects scales for progressive image compression~\cite{arpc}. Ada-TokenCom jointly adapts prefix length and the modulation and coding scheme (MCS) under a long-term symbol budget~\cite{adatokencom}, and ResARC codes residuals to compensate for suffix-generation errors~\cite{resarc}. Deep joint source--channel coding (JSCC) methods such as SwinJSCC map images directly to channel symbols and serve as system-level baselines~\cite{swinjscc}. For a fixed digital frame, the open question is how to terminate a decodable prefix within a scale and use the recovered boundary tokens at the receiver.

This letter proposes NeST-Com. Its frame-fitting rule uses a coder-specific nondecreasing-length property to select the longest independently terminated within-scale prefix whose arithmetic-code length fits the protected frame body. The receiver incorporates the recovered boundary tokens through VAR's known-token overwrite to correct residuals and condition finer-scale predictions. Experiments on 500 ImageNet images and matched controls examine how capacity utilization and boundary-token use contribute to the gains and quantify the quality--computation tradeoff.
